\documentclass[11pt]{article}
\usepackage{enumitem}
\usepackage{listings}
\title{Computer Abstraction and Technology.}
\lstset{basicstyle=\ttfamily, breaklines=true}
\begin{document}
\maketitle
\begin{enumerate}
    \item [1.1]\textbf{[2] $<$§1.1$>$ Aside from the smart cell phones used by a billion people, list and describe four other types of computers.}
    \begin{enumerate}
        \item Personal computers.
        {\fontfamily{cmtt}\selectfont They emphasize on good perfomance to a single user at a low cost. They usually execute third party software.}
        \item Embedded computers.
        {\fontfamily{cmtt}\selectfont The largest classes of computers. They are designed to run one application that's usually integrated with the hardware
        and delivered as a single system.}
        \item Servers 
        {\fontfamily{cmtt}\selectfont Usually accessed via network. Servers are oriented towards carrying a sizable complex workload which may be a single application that's 
        encompasses a scientific application or an engineering application - or it may handle multiple small jobs. Servers are built from the same technology as destop computers
        but provide for greater conputing, storage, and input/output capacity. Servers place more emphasies on dependability. Servers span the widest range in terms of cost and capability.}
        \item Supercomputers.
        {\fontfamily{cmtt}\selectfont Consist of thousands to millions of processors and terabytes of storage. They are used for high-end scientific and engineering calculations,
        such as weather forecasting, oil exploration, protein structure determination and other large scale problems.}
    \end{enumerate}
    \item [1.2]\textbf{[5] $<$§1.2$>$ The eight great ideas in computer architecture are similar to ideas from other fields. Match the eight ideas from computer architecture, “Design for 
    Moore's Law,” “Use Abstraction to Simplify Design,” “Make the Common Case Fast,” “Performance via Parallelism,” “Performance via Pipelining,” “Performance 
    via Prediction,” “Hierarchy of Memories,” and “Dependability via Redundancy” to the following ideas from other fields:}
    \begin{enumerate}
        \item Assembly lines in automobile manufacturing\\
        {\fontfamily{cmtt}\selectfont Perfomance via Pipelining}
        \item Suspension bridge cables\\
        {\fontfamily{cmtt}\selectfont Perfomance via Parallelism}
        \item Aircraft and marine navigation systems that incorporate wind information\\
        {\fontfamily{cmtt}\selectfont Use Abstraction to Simplify Design}
        \item Express elevators in buildings\\
        {\fontfamily{cmtt}\selectfont Make the Common Case Fast}
        \item Library reserve desk\\
        {\fontfamily{cmtt}\selectfont Hierarchy of Memories}
        \item Increasing the gate area on a CMOS transistor to decrease its switching time\\
        {\fontfamily{cmtt}\selectfont Design for Moore's Law}
        \item Adding electromagnetic aircraft catapults (which are electrically powered
              as opposed to current steam-powered models), allowed by the increased power
              generation offered by the new reactor technology.\\
        {\fontfamily{cmtt}\selectfont Performance via Prediction}
        \item Building self-driving cars whose control systems partially rely on existing sensor
              systems already installed into the base vehicle, such as lane departure systems and
              smart cruise control systems.\\
        {\fontfamily{cmtt}\selectfont Dependability via Redundancy}
    \end{enumerate}
    \item[1.3]\textbf{[2]$<$§1.3$>$ Describe the steps that transform a program written in a high-level language 
                                 such as C into a representation that is directly executed by a computer processor.}
        \begin{enumerate}
            \item {\fontfamily{cmtt}\selectfont High-level language -$>$ Compiled to -$>$ Assembly language -$>$ Assembler(Assembles the Assembly language to) -$>$ Machine language(Binary instructions)}
        \end{enumerate}
\end{enumerate}
\end{document}